\begin{statementbox}

	\textbf{\textcolor{CStatement}{条件}}
	\begin{description}[style=nextline, leftmargin=2.8em, labelsep=0.5em, itemsep=0.35em]
		\item[\textbf{\textcolor{CStatement}{条件 1（命题形式）：}}] 原命题为“若 $p$，则 $q$”形式，其中 $p$ 是条件，$q$ 是结论。通常命题隐含全称量词，即 $\forall x\,(p(x)\to q(x))$。
	\end{description}

	\textbf{\textcolor{CStatement}{结论}}
	\begin{description}[style=nextline, leftmargin=2.8em, labelsep=0.5em, itemsep=0.45em]
		\item[\textbf{\textcolor{CStatement}{结论一（命题否定）：}}]
			\textbf{\textcolor{CStatement}{直观描述：}}
			否定一个命题只需举一个反例，即条件成立但结论不成立。
			\[
				\lnot(p\to q)\equiv p\land\lnot q
			\]
			若原命题隐含全称量词，则其否定为存在量词形式：
			\[
				\lnot(\forall x\,(p(x)\to q(x)))\equiv \exists x\,(p(x)\land\lnot q(x))
			\]

		\item[\textbf{\textcolor{CStatement}{结论二（否命题）：}}]
			\textbf{\textcolor{CStatement}{直观描述：}}
			否命题是原命题的条件和结论分别取否定后组成的新命题。
			\[
				\lnot p\to\lnot q
			\]
			注：否命题与原命题不一定同真假。

		\item[\textbf{\textcolor{CStatement}{推广结论（逆否等价）：}}]
			\textbf{\textcolor{CStatement}{直观描述：}}
			原命题与其逆否命题始终同真同假，是最重要的逻辑等价关系之一。
			\[
				(p\to q)\equiv(\lnot q\to\lnot p)
			\]
	\end{description}

\end{statementbox}
